% STATUS: VERIFIED / AUDITED
% AUDIT FILE: audits/A05_q2_mode_selection.tex
% ORIGINAL AUDIT VERDICT: CORE MINOR REVISION; REALIZABILITY VERIFIED;
%   OVERALL MINOR REVISION
% AUDIT CLOSURE: A05 sole required issue E01 applied verbatim
% FINAL DISPOSITION: VERIFIED / AUDITED AFTER MINOR REVISION
% ORIGINAL AUDITED VERSION: P05 v1.0; 2026-10-02; 26848 bytes; 591 lines
% ORIGINAL AUDITED FULL-FILE SHA-256: b7123736242554bd7fc1f6a56d8788b6c5413a25930e5360d7ea5babe44184d4
% ORIGINAL AUDITED BODY SHA-256: 25ab3e4ed5d871ab2884210e6d8e7f3631acee3f7a1c662d3c33de441048810d
% CURRENT VERSION: P05 v1.1; 2026-10-02
% CURRENT BODY SHA-256: 7e1fd7765ed81e9c5f24ebe07e450e3fd76caf6b16264aac3c3a21edacd8bb97
% CURRENT HASH: full-file SHA-256 recorded in project_manifest.md and
%   the closure report to avoid a self-referential full-file hash.
% BODY HASH SCOPE: bytes from the first line-anchored section command to EOF.
% CLOSURE DATE: 2026-10-02
% PRESERVATION: the sole body change is E01 (candidate versus selected labels).
% STATUS NOTE: historical DRAFT/audit-pending body prose is preserved;
%   this metadata and the current manifest govern status.
% PROVENANCE NOTE: A05 U01 is nonblocking. Original validation companion
%   unavailable to the auditor; historical run claims not independently reproduced.
%   Closure neither reruns historical validation nor claims to resolve U01.
% DOWNSTREAM: P08 prerequisite cleared; no later proof module begun.
% DEPENDS ON: P01-P04 VERIFIED / AUDITED
% MODEL: canonical demand-modulated model, q=2
% LITERATURE DEPENDENCIES: frozen Phase 1B/1C novelty audit only
% NOVELTY ROLE: MAIN THEOREM TARGET only in strengthened realizability form
% NUMERICAL DEPENDENCIES: none
% NEXT DEPENDENCIES: P08, P07, P11
% VERSION: P05 v1.1; 2026-10-02; A05 E01 corrected
% CLAIMS: C045-C048; P05-D1-P05-D5 in claim_registry.md
% SCOPE: fixed positive synchronous equilibrium, exact finite n, all ties,
% canonical realizability and local switch geometry; no first-Hopf theorem.

\section{P05: Exact step-two mode selection}
\label{sec:P05}
\noindent\textbf{Status: DRAFT; independent audit pending.}
Fix the canonical demand-modulated model with $q=2$, $n\ge5$ and
$a,k,c,g,h>0$. Let $x^*\one$ be a specified positive synchronous
equilibrium and put $s=1+x^*>1$. Import only the audited P04 coefficients
and real dispersion relation:
\begin{equation}
\begin{gathered}
 C=\frac cs>0,\qquad G=\frac{gx^*}{s^2}>0,\qquad H=hx^*>0,\\
 r_m=d+C\cos\theta_m+(G-H)\cos(2\theta_m),\qquad
 \theta_m=\frac{2\pi m}{n}.
\end{gathered}
 \label{eq:P05-input}
\end{equation}
Here $r_m=\Ree\lambda_m$ and $d$ is P04's diagonal coefficient.
All selection conclusions below are derived from \eqref{eq:P05-input}.
The maximization is over all Fourier labels, including mode zero, and
selection means the complete set maximizing the real part at this
equilibrium, retaining all ties. No coprimality condition is used.
P01 supplies smoothness, P02 supplies the equilibrium relation, and P03--P04
supply the homogeneous coefficients. No later proof module is used.

\paragraph{Notation and novelty scope.}
Use $T=G-H$ and the dimensionless scalar
\begin{equation}
 \chi=\frac{4(H-G)}{C}=-\frac{4T}{C}.
 \label{eq:P05-control}
\end{equation}
The letter $D$ remains available for P02's discriminant, and $\rho$
retains the group-action meaning in Section~\ref{sec:notation}.
Here $j$ denotes a Fourier representative, rather than a node index.
The frozen Phase 1B/1C verdict treats quadratic maximization as routine.
The intended main-theorem content is the complete finite-grid classification,
its realization by positive canonical parameters, and local switch geometry.
No priority claim or new literature assessment is made.

\subsection{Quadratic dispersion relation}
\label{slot:P05-1}
Set $u_m=\cos\theta_m$. The identity $\cos(2\theta)=2\cos^2\theta-1$
gives
\begin{equation}
 r_m=R(u_m),\qquad
 \boxed{R(u)=d-T+Cu+2Tu^2.}
 \label{eq:P05-quadratic}
\end{equation}
The additive constant is $d-T$; the part controlling selection is
$Cu+2Tu^2$. In particular, differences of growth rates and the complete
maximizer set depend only on $C,T$, not on $d$. This is an algebraic
selection statement, not permission to vary $d$ independently in the model.

Put $N=\lfloor n/2\rfloor$ and $u_j=\cos(2\pi j/n)$ for $0\le j\le N$.
Since cosine is strictly decreasing on $[0,\pi]$,
\begin{equation}
 1=u_0>u_1>\cdots>u_N,\qquad
 u_N=\begin{cases}-1,&n\text{ even},\\-\cos(\pi/n),&n\text{ odd}.
 \end{cases}
 \label{eq:P05-grid}
\end{equation}
Equation~\eqref{eq:P05-input} gives $r_{n-m}=r_m$ for $1\le m<n$.
Define the distinct label sets
\begin{equation}
 L_j=\begin{cases}
 \{0\},&j=0,\\
 \{n/2\},&n\text{ even and }j=N,\\
 \{j,n-j\},&\text{otherwise}.
 \end{cases}
 \label{eq:P05-labels}
\end{equation}
For a maximizing representative set $S$, the full dominant-label set is
$\bigcup_{j\in S}L_j$, a disjoint union.

\subsection{Continuous maximizer}

If $T\ge0$, the polynomial is convex (linear when $T=0$). Explicitly,
\[
 R(1)-R(-1)=2C>0,
 \qquad R(1)-R(u)=(1-u)\bigl[C+2T(1+u)\bigr]>0
 \quad(-1\le u<1).
\]
Thus $j=0$ is the unique dominant representative, including $T=0$ and
odd $n$, when $-1$ itself is absent from the grid.

If $T<0$, then $H-G>0$, $R''=-4(H-G)<0$, and
\begin{equation}
 R'(u)=C-4(H-G)u,\qquad
 u_*:=\frac{C}{4(H-G)}=\frac1\chi>0.
 \label{eq:P05-vertex}
\end{equation}
If $u_*\ge1$, the derivative is positive on $[-1,1)$, hence $j=0$
is still uniquely selected. The boundary $u_*=1$ is exactly
$H-G=C/4$, or $\chi=1$; the zero derivative at $u=1$ does not give
a second grid maximizer. If $0<u_*<1$, the continuous maximizer on
$[-1,1]$ is the interior vertex. A discrete change need not occur
when this vertex first leaves the endpoint.

\subsection{Exact discrete selection rule}
\label{slot:P05-2}
For $T<0$, direct completion of the square gives
\begin{equation}
 R(u)=d+(H-G)+\frac{C^2}{8(H-G)}
           -2(H-G)(u-u_*)^2.
 \label{eq:P05-square}
\end{equation}
Consequently maximizing the growth rate is exactly minimizing
$|u_j-u_*|$. This holds for every $u_*>0$, including $u_*\ge1$.
Distances are in cosine space; the grid in that space is nonuniform.
No angular rounding rule is invoked.

Define the strictly decreasing midpoint sequence
\begin{equation}
 \beta_j=\frac{u_j+u_{j+1}}2
 =\cos\frac\pi n\cos\frac{(2j+1)\pi}{n},\qquad 0\le j<N.
 \label{eq:P05-midpoints}
\end{equation}
Strict decrease follows from
$\beta_j-\beta_{j+1}=(u_j-u_{j+2})/2>0$.

\begin{theorem}[A: Exact selection, including endpoints and all ties]
\label{thm:P05-selection}
Under the stated canonical equilibrium assumptions, if $T\ge0$ then
$S=\{0\}$. If $T<0$, the selected representatives are exactly those
whose closed cells below contain $u_*$:
\begin{align}
 V_0&=[\beta_0,\infty),\label{eq:P05-cells}\\
 V_j&=[\beta_j,\beta_{j-1}]\quad(1\le j<N),\nonumber\\
 V_N&=(-\infty,\beta_{N-1}].\nonumber
\end{align}
A representative is the unique maximizer precisely when $u_*$ lies in
the corresponding open cell. At $u_*=\beta_j$ the selected set is
exactly $\{j,j+1\}$. No nonadjacent or triple representative tie can
be maximal. Only positive midpoints can be attained when $T<0$.
The full Fourier-label set is given by \eqref{eq:P05-labels}.
\end{theorem}
\begin{proof}
The preceding endpoint comparison proves the $T\ge0$ assertion.
For $T<0$ use \eqref{eq:P05-square}. For any adjacent grid points,
\[
 (u_j-u_*)^2-(u_{j+1}-u_*)^2
       =2(u_j-u_{j+1})(\beta_j-u_*).
\]
Since $u_j>u_{j+1}$, representative $j$ is closer exactly when
$u_*>\beta_j$, with equality exactly at the midpoint.
For $1\le j<N$ the two neighboring comparisons require
$\beta_j\le u_*\le\beta_{j-1}$. These inequalities imply
$u_{j+1}<u_*<u_{j-1}$; every grid point beyond either neighbor is
even farther in that direction. Thus the neighboring comparisons are
also sufficient. For $j=0$, the single inequality $u_*\ge\beta_0$
makes every $u_\ell$, $\ell>1$, farther than $u_1$; the analogous
argument at $N$ uses $u_*\le\beta_{N-1}$. Strict inequalities give
uniqueness, and equality gives exactly the stated adjacent pair.

For completeness, if two nonadjacent representatives had equal maximal
value, a grid point strictly between their cosines would have strictly
larger value by strict concavity. This is impossible. Three maximizers
would include a nonadjacent pair. Equality of two nonadjacent values
can occur away from the maximum; the theorem excludes only maximal ties.
Finally $u_*>0$ excludes all nonpositive midpoint equalities.
\end{proof}

\subsection{Mode ties and selection cells}
\label{slot:P05-3}
The real-part comparison also has the useful exact form
\begin{equation}
 R(u_j)-R(u_{j+1})
       =C(u_j-u_{j+1})(1-\chi\beta_j).
 \label{eq:P05-adjacent-gap}
\end{equation}
For $T<0$ an adjacent \emph{maximal} tie therefore occurs if and only
if $u_*=\beta_j>0$, equivalently $\chi=1/\beta_j$.
The word maximal is essential: adjacent equalities outside the concave
regime need not be ties at the spectral abscissa.

Let $w_j=|L_j|$. A strict cell has $w_j$ selected Fourier labels; a
representative tie has $w_j+w_{j+1}$. Thus two nonsingleton
representatives give four labels, a tie involving $0$ gives three,
and an even-$n$ endpoint would contribute only one label. Below we prove
that the $N=n/2$ endpoint is never selected here. In the canonical
reachable cells every nonzero selected representative gives exactly two
labels, and every reachable nonzero/nonzero switch gives four.
These counts are counts of labels, not of distinct numerical eigenvalues.

\subsection{Reachable Fourier representatives}

\begin{theorem}[B: Reachable cells and parity]
\label{thm:P05-reachable}
The full list of representatives whose nearest-cosine cell intersects
$(0,1)$, and also the full list attainable at some canonical positive
synchronous equilibrium, is
\begin{equation}
 \boxed{0,1,\ldots,J,\qquad J=j_{\max}(n)=\left\lfloor\frac{n+1}{4}\right\rfloor.}
 \label{eq:P05-jmax}
\end{equation}
Each is attainable as the unique representative. Precisely the midpoints
$\beta_0,\ldots,\beta_{J-1}$ are positive and give attainable switches.
Every $j>J$, in particular the even-$n$ singleton $n/2$, is excluded.
Canonical attainability and robustness are proved in Theorem~\ref{thm:P05-realization}.
\end{theorem}
\begin{proof}
For $1\le j\le N$, the cell intersects $(0,1)$ if and only if
$\beta_{j-1}>0$. Necessity follows from its upper endpoint. Sufficiency
follows by taking a point strictly between
$\max(0,\beta_j)$ and $\beta_{j-1}$ when $j<N$, with the endpoint
case interpreted using $V_N$. The upper endpoint is less than $1$.
The cell for $0$ intersects $(0,1)$ since $0<\beta_0<1$.
By \eqref{eq:P05-midpoints} and $\cos(\pi/n)>0$,
\[
 \beta_j>0\ \Longleftrightarrow\ 2j+1<n/2
 \ \Longleftrightarrow\ j<(n-2)/4.
\]
Indeed the relevant midpoint angles lie in $(0,\pi)$, where cosine
changes sign only at $\pi/2$. Hence
$\beta_{j-1}>0$ is equivalent to $j<(n+2)/4$, giving exactly
\eqref{eq:P05-jmax}. In particular $J<N$ for every $n\ge5$,
$\beta_{J-1}>0$, and $\beta_J\le0$. All retained cell intersections
have nonempty interiors; there is no additional representative reachable
only through an isolated tie. When $T\ge0$ or $u_*\ge1$ only $0$
is selected, so neither regime enlarges the list.
\end{proof}

The parity information, with $n=4\ell+r$ in each row, is exact:
\begin{center}
\renewcommand{\arraystretch}{1.2}
\begin{tabular}{@{}c c c c@{}}
\toprule
$r$ & $J$ & $u_J$ & Nearest representatives at the formal point $u_*=0$\\
\midrule
$0$ & $\ell$ & $0$ & $\ell$\\
$1$ & $\ell$ & $\sin(\pi/(2n))>0$ & $\ell$\\
$2$ & $\ell$ & $\sin(\pi/n)>0$ & $\ell,\ell+1$\\
$3$ & $\ell+1$ & $-\sin(\pi/(2n))<0$ & $\ell+1$\\
\bottomrule
\end{tabular}
\end{center}
To verify the last column independently, the two cosines straddling zero
in the odd cases are
\[
 \begin{array}{ll}
 n=4\ell+1:&u_\ell=\sin(\pi/(2n)),\quad
             u_{\ell+1}=-\sin(3\pi/(2n)),\\
 n=4\ell+3:&u_\ell=\sin(3\pi/(2n)),\quad
             u_{\ell+1}=-\sin(\pi/(2n)).
 \end{array}
\]
All their arguments are in $(0,\pi/2)$, where sine increases strictly.
All remaining cosines are farther from zero by monotonicity of the grid.
For $4\mid n$ there is an exact zero cosine at $j=n/4$, and its cell
contains a positive interval. This representative is a conjugate pair,
not the real singleton $n/2$. For $n=4\ell+2$, the straddling values
are $\pm\sin(\pi/n)$, and $\beta_\ell=0$.

\paragraph{Large restraint at fixed finite dimension.}
For fixed $C>0$, as $H-G\to+\infty$, the vertex decreases to zero
and $J$ is the unique dominant representative for all sufficiently large
finite $H-G$, in every parity class. The formal zero-vertex tie for
$n=4\ell+2$ is never attained with finite positive $C,H-G$.
In fact the two competing real parts have the strictly positive,
$H-G$-independent difference
\[
 r_\ell-r_{\ell+1}=2C\sin(\pi/n)>0.
\]
Their squared cosines are equal, so this follows directly from
\eqref{eq:P05-quadratic}. Only a leading-order normalization by $H-G$
loses that difference. Thus a formal limiting nearest-grid tie must not
be added to the finite-parameter dominant set. No $n\to\infty$ limit is taken.
For $n\equiv3\pmod4$, $J$ is the sole reachable representative with a
negative cosine; for the other residue classes no negative-cosine
representative is reachable.

\paragraph{Complete regions in one dimensionless parameter.}
Put $\alpha_j=1/\beta_j$ for $0\le j<J$. Then
$1<\alpha_0<\cdots<\alpha_{J-1}$, and for \emph{all} real $\chi$
allowed by the canonical coefficients the entire selection rule is
\begin{equation}
 \begin{array}{c|c}
 \text{Condition}&S\\\hline
 \chi<\alpha_0&\{0\}\\
 \alpha_{j-1}<\chi<\alpha_j\quad(1\le j<J)&\{j\}\\
 \chi>\alpha_{J-1}&\{J\}\\
 \chi=\alpha_j\quad(0\le j<J)&\{j,j+1\}.
 \end{array}
 \label{eq:P05-regions}
\end{equation}
The middle row is empty when $J=1$. Negative and zero $\chi$ are covered
by the first row. For positive $\chi$, the inequalities follow by inverting
positive midpoint bounds, reversing their order. Thus the continuous
vertex boundary $\chi=1$ is strictly inside the uniform selection region;
the first discrete switch is at $\alpha_0>1$.

\subsection{Canonical parameter realization}

The concave condition and vertex in original variables are
\begin{equation}
 T<0\ \Longleftrightarrow\ h>\frac{g}{s^2},\qquad
 u_*=\frac{c/s}{4x^*(h-g/s^2)}
     =\frac{cs}{4x^*(hs^2-g)},\qquad
 \chi=\frac{4x^*(hs^2-g)}{cs}.
 \label{eq:P05-canonical-control}
\end{equation}
All denominators in the vertex formula are strictly positive in this
regime. For $0\le j<J$, the switch equations are exactly
\begin{equation}
 H-G=\frac{C}{4\beta_j},\qquad
 \boxed{h=\frac{g}{(1+x^*)^2}
       +\frac{c}{4(1+x^*)x^*\beta_j}.}
 \label{eq:P05-switch}
\end{equation}
They are subject to P02's equilibrium condition
\begin{equation}
 a-(k+h)x^*+\frac{c-g}{1+x^*}=0.
 \label{eq:P05-equilibrium}
\end{equation}
For $\beta_j\le0$ there is no switch in the $T<0$ regime: division
by zero is not meaningful and negative $\beta_j$ would require
$H-G<0$. Such negative-midpoint equalities in the convex regime, if
realizable as equal lower values, are not dominant-mode switches.

\begin{theorem}[C: Canonical realizability, open cells and switch geometry]
\label{thm:P05-realization}
Every strict cell and every switch in \eqref{eq:P05-regions} is nonempty
in the canonical positive-equilibrium parameter space. Each strict cell
has a realization persisting on a nonempty open set of positive physical
parameters $(a,k,c,g,h)$ with a specified smooth positive synchronous
branch. Both $T>0$ and $T<0$, and the boundary $T=0$, are realizable.
Every positive-midpoint switch is locally a smooth codimension-one
hypersurface in physical parameter space along its specified branch;
all its positive-equilibrium points are regular branch points. No
nonpositive-midpoint dominant switch is realizable. These are local
branch statements and assert no global identification of distinct branches.
\end{theorem}
\begin{proof}[Proof of realizability]
Choose any $x>0$, any $c,g>0$, and any desired vertex $v>0$. Define
\begin{equation}
 s=1+x,\qquad h=\frac{g}{s^2}+\frac{c}{4sxv},\qquad
 k>\max\left\{0,\frac{c-g}{sx}-h\right\},\qquad
 a=(k+h)x-\frac{c-g}{s}.
 \label{eq:P05-construction}
\end{equation}
Then $h,k,a$ are strictly positive: the displayed strict lower bound on
$k$ is exactly what is required for $a>0$, together with $k>0$.
There is no upper bound on $k$, so this choice is always possible.
Equation~\eqref{eq:P05-equilibrium} vanishes identically at $x^*=x$;
substitution into the original vector field gives, in every component,
\[
 ax-kx^2+\frac{cx}{s}-\frac{gx}{s}-hx^2
 =x\left[a-(k+h)x+\frac{c-g}{s}\right]=0.
\]
Also $H-G=c/(4sv)>0$ and $u_*=v$. Hence every retained open
Voronoi cell and every positive midpoint can be prescribed exactly.
The construction does not treat $C,G,H$ as independent of an equilibrium.

A particularly simple exact family, used below, is
\begin{equation}
 x^*=1,\qquad c=g=4,\qquad h=1+\chi/2,\qquad
 k>0,\qquad a=k+1+\chi/2,\qquad \chi>-2.
 \label{eq:P05-simple-family}
\end{equation}
Here $C=2$, $G=1$, $H=1+\chi/2$, so the dimensionless parameter
is exactly the displayed $\chi$. Every positive $\chi$ is realized.
At $k=1$, choosing $\chi=-1,0,1$ realizes respectively $T>0$,
$T=0$, and $T<0$ with $u_*=1$, all selecting $0$.
Thus the uniform region has examples in both sign regimes.
The geometry and openness assertions are proved next.
\end{proof}

\subsection{Realizable switch hypersurfaces}

Use the equilibrium coordinates $(x,k,c,g,h)$ on the open set
\begin{equation}
 \mathcal U_{\rm eq}=
 \left\{(x,k,c,g,h)\in(0,\infty)^5:
     (k+h)x-\frac{c-g}{1+x}>0\right\}.
 \label{eq:P05-equilibrium-chart}
\end{equation}
The map
\[
 \Phi(x,k,c,g,h)=\left((k+h)x-\frac{c-g}{1+x},\ k,c,g,h\right)
\]
projects this parameterization of the positive equilibrium graph onto
physical parameters. With $s=1+x$ its Jacobian determinant is
\begin{equation}
 \det D\Phi=k+h+\frac{c-g}{s^2}.
 \label{eq:P05-chart-determinant}
\end{equation}
At every point with $H-G>0$, this is
$k+(h-g/s^2)+c/s^2>0$. Equivalently, for the scalar equilibrium
function $B(x)=a-(k+h)x+(c-g)/(1+x)$,
$B_x=-\det D\Phi<0$. The inverse function theorem makes $\Phi$ a
local diffeomorphism there, and the implicit function theorem gives a
specified smooth positive branch $x=x^*(a,k,c,g,h)$. This regularity
uses only the scalar relation, not full-system invertibility or a Hopf result.

Strict inequalities in \eqref{eq:P05-regions} define open subsets of
$\mathcal U_{\rm eq}$. Every strict cell has a $\chi>0$ point from
\eqref{eq:P05-construction}, including the uniform cell, so local
projection supplies a nonempty open set of physical parameters for each.
In fact every concave-regime point of a strict cell has this local property.
For a $T>0$ uniform example, the family at $k=1,\chi=-1$ has
$\det D\Phi=k+h=3/2>0$; strict $T>0$ persists in a neighborhood.
No global openness assertion at an arbitrary scalar fold in the uniform
region is required or inferred.

For $\beta_j>0$ define in equilibrium coordinates
\begin{equation}
 \Psi_j(x,k,c,g,h)=h-\frac{g}{(1+x)^2}
                  -\frac{c}{4(1+x)x\beta_j}.
 \label{eq:P05-level-set}
\end{equation}
Its partial derivative with respect to $h$, holding $x,k,c,g$ fixed,
is exactly $1$. Hence its zero set is a smooth codimension-one regular
level set within $\mathcal U_{\rm eq}$. It is nonempty by the construction,
and every one of its points has $H-G>0$, so the local diffeomorphism
$\Phi$ transfers this conclusion to physical parameter space on the
chosen branch. Moreover, at fixed $x,c,g,k$,
\[
 \frac{\partial\chi}{\partial h}=\frac{4(1+x)x}{c}>0,
 \qquad
 \frac{\partial}{\partial h}(r_j-r_{j+1})
       =-4x(u_j-u_{j+1})\beta_j<0.
\]
Here $a$ varies according to \eqref{eq:P05-equilibrium}; this is a
legitimate direction in physical parameter space. Distinct midpoint
thresholds and the strict gap to all other representatives at the tie
show that the two sides locally select $j$ and $j+1$, respectively.
This proves genuine local changes of dominance, and finishes
Theorem~\ref{thm:P05-realization}. No global manifold structure after
forgetting branch labels is claimed.

\subsection{Exact examples and small-dimension checks}

In \eqref{eq:P05-simple-family} take $k=1$. The tuple
$(a,k,c,g,h)=(2,1,4,4,1)$ selects $0$ in every dimension.
For $n=8$, $(3,1,4,4,2)$ has $\chi=2$, $u_*=1/2$, and selects
representative $1$, since
$\sqrt2/4<1/2<(2+\sqrt2)/4$.
For $n=6$, $(8/3,1,4,4,5/3)$ has $\chi=4/3=1/\beta_0$ and
selects exactly labels $\{0,1,5\}$.
For $n=8$, $(2+\sqrt2,1,4,4,1+\sqrt2)$ has
$\chi=2\sqrt2=1/\beta_1$ and selects $\{1,2,6,7\}$.
Every example has $x^*=1$, all five coefficients positive, and original
DM residual $a-k-h+(c-g)/2=0$.

For exact highest-mode examples across all four residue classes, use
$k=1$ in the same family with $\chi=2/\beta_{J-1}$. Then
$u_*=\beta_{J-1}/2\in(0,\beta_{J-1})$, while $\beta_J\le0$,
so $J$ is strict. In particular the following are exact specifications:
\begin{center}
\begin{tabular}{@{}c c c c@{}}
\toprule
$n$ & $J$ & $\chi$ & $u_J$\\\midrule
$8$ & $2$ & $4\sqrt2$ & $0$\\
$9$ & $2$ & $4/\cos(\pi/9)$ & $\sin(\pi/18)$\\
$10$ & $2$ & $2/[\cos(\pi/10)\cos(3\pi/10)]$ & $\sin(\pi/10)$\\
$11$ & $3$ & $2/[\cos(\pi/11)\cos(5\pi/11)]$ & $-\sin(\pi/22)$\\
\bottomrule
\end{tabular}
\end{center}
The corresponding tuples are exactly
$(2+\chi/2,1,4,4,1+\chi/2)$. No numerical search is used.

The next table records the complete representative grid and positive
midpoints for $5\le n\le12$. Within each row only, write
$\gamma_r=\cos(r\pi/n)$. All listed midpoint values are exact.
Their reciprocal values are the switch thresholds in
\eqref{eq:P05-regions}; the reachable set is $\{0,\ldots,J\}$.
The terminal midpoint $\beta_J$ is negative except at $n=6,10$,
where it is zero and unreachable. The five-node row is solely a check;
the golden-ratio specialization is reserved for P07.
\begin{center}
\small
\renewcommand{\arraystretch}{1.3}
\begin{tabular}{@{}c l c l@{}}
\toprule
$n$ & $(u_0,\ldots,u_N)$ & $J$ & $(\beta_0,\ldots,\beta_{J-1})$\\
\midrule
$5$ & $(1,\gamma_2,-\gamma_1)$ & $1$ & $(\gamma_1^2)$\\
$6$ & $(1,1/2,-1/2,-1)$ & $1$ & $(3/4)$\\
$7$ & $(1,\gamma_2,-\gamma_3,-\gamma_1)$ & $2$ & $(\gamma_1^2,\gamma_1\gamma_3)$\\
$8$ & $(1,\sqrt2/2,0,-\sqrt2/2,-1)$ & $2$ & $((2+\sqrt2)/4,\sqrt2/4)$\\
$9$ & $(1,\gamma_2,\gamma_4,-1/2,-\gamma_1)$ & $2$ & $(\gamma_1^2,\gamma_1/2)$\\
$10$ & $(1,\gamma_2,\gamma_4,-\gamma_4,-\gamma_2,-1)$ & $2$ & $(\gamma_1^2,\gamma_1\gamma_3)$\\
$11$ & $(1,\gamma_2,\gamma_4,-\gamma_5,-\gamma_3,-\gamma_1)$ & $3$ & $(\gamma_1^2,\gamma_1\gamma_3,\gamma_1\gamma_5)$\\
$12$ & $(1,\sqrt3/2,1/2,0,-1/2,-\sqrt3/2,-1)$ & $3$ & $((2+\sqrt3)/4,(\sqrt3+1)/4,1/4)$\\
\bottomrule
\end{tabular}
\end{center}

\subsection{Dominant modes versus first instability}

The preceding results concern $\argmax_m r_m$ at a fixed equilibrium.
The scalar $d$ cancels from that comparison but affects whether the
maximum is negative, zero, or positive. A first loss of stability would
require a specified oriented physical path, a positive equilibrium branch,
a prior interval with $\max_m r_m<0$, and a crossing to a positive
maximum. Neither a selection switch nor a real-part tie establishes that.

Indeed every reachable switch can occur on a path that stays strictly
stable. Fix its threshold $\alpha_j$, and choose a compact interval
$[b_-,b_+]\subset(0,\infty)$ with $\alpha_j$ in its interior and no
other threshold in it. In \eqref{eq:P05-simple-family} let $\chi$
increase across this interval, hold $k>b_+/2$ fixed, and keep $x^*=1$.
The equilibrium is regular, all parameters remain positive, and P04's
verified equilibrium diagonal gives $d=-k-3$. Thus directly from
\eqref{eq:P05-input},
\begin{equation}
 r_m=-k-3+2\cos\theta_m-\frac\chi2\cos(2\theta_m)
       \le-k-1+\frac{b_+}{2}<-1.
 \label{eq:P05-stable-switch}
\end{equation}
The dominant representative changes from $j$ to $j+1$, with the exact
tie at $\alpha_j$, while every mode remains strictly stable. This also
shows that every strict cell has a stable canonical realization, since
$k$ can be enlarged without changing $\chi$ or $x^*=1$ in this family.
No first-instability subresult is asserted; that separate crossing question
is deferred to P11/P12 under their branch/path gates.

\paragraph{Real-part ties are not eigenvalue coincidences.}
At every switch P05 proves $r_j=r_{j+1}$ only. P04 supplies the signed
imaginary parts
\[
 \omega_m=-C\sin\theta_m+(G+H)\sin(2\theta_m).
\]
For particular labels, equality of the complex eigenvalues would additionally
require equality of these signed imaginary parts. For conjugate labels the
sign changes. None of the selection, realization, or geometry proofs needs
an imaginary-part condition. Numerical eigenvalue simplicity, nonzero
critical frequency, simultaneous critical pairs and Hopf exclusions remain
P11's tasks. Counting two or four selected labels does not certify a simple
pair or a multiple-eigenvalue collision.

\subsection{Checks and possible failure points}

\paragraph{Internal analytic checks.}
Expansion of \eqref{eq:P05-square}, its derivative, and the explicit
endpoint differences check the quadratic and continuous regimes.
The squared-distance difference checks the orientation of every midpoint
inequality; strict concavity excludes nonadjacent maximal ties.
The midpoint product in \eqref{eq:P05-midpoints} and the four straddling
cosine calculations independently check $J$ and the negative-cosine case.
Equations~\eqref{eq:P05-canonical-control}--\eqref{eq:P05-construction}
check the canonical substitutions, positivity and original vector-field
residual. The determinant and $\partial_h\Psi_j=1$ check the branch
regularity and codimension statements. Singleton bookkeeping uses actual
Fourier labels; the eigenvalue discussion does not infer simplicity.

\paragraph{Reproducible corroboration.}
The separately delivered P05 validation support contains a deterministic
checker, exact rational symbolic identities, high-precision all-label
comparisons for every $5\le n\le40$, canonical construction residuals,
and build/provenance records. It samples several points in every strict
cell, both sides of every switch and the switches themselves, as well as
$T=0$, $u_*=1$, near-zero vertices and every residue class modulo four.
Tie inputs are defined by exact trigonometric expressions; numerical
comparisons at those inputs are still finite-precision corroboration.
The analytic proofs above do not depend on computed values. Run totals,
tolerances and environment are recorded in the support output and manifest.
Compilation and these internal checks do not verify or independently audit P05.

\paragraph{Priority checks for the independent P05 auditor.}
\begin{enumerate}[label=(\arabic*),leftmargin=*]
\item Authenticate the current P01--P04 closure versions, preserve all their
bodies, and keep this P05 version DRAFT until independently audited.
\item Check the signs in the quadratic, square completion and adjacent gap;
retain mode zero, $T=0$, $u_*=1$, and the actual discrete first threshold.
\item Check the descending-grid Voronoi inequalities, both endpoint cells,
all maximal ties, and the exclusion of nonadjacent and triple maxima.
\item Derive $J=\lfloor(n+1)/4\rfloor$ independently. Inspect the negative
cosine for $n\equiv3\pmod4$, zero cosine for $4\mid n$, and the unattainable
zero-midpoint tie for $n\equiv2\pmod4$, including its persistent finite gap.
\item Verify every constructed coefficient is positive and substitute into
the original DM equation, rather than treating spectral coefficients as free.
\item Distinguish openness in equilibrium coordinates from openness after
projection. Check $\det D\Phi>0$ throughout the concave regime and
$\partial_h\Psi_j=1$ on each switch; assert only the stated local geometry.
\item Count singleton and paired labels correctly; real-part ties do not
establish eigenvalue equality, nonzero frequency, or a simple Hopf pair.
\item Check the stable switch path and ensure no hidden dependence on
P07, P08 or P11--P13, no first-instability assertion, and no novelty claim
for the bare quadratic maximization.
\end{enumerate}
